\section{Optical cofactors, multiple marks and simultaneous endpoints}
\label{sec:v110-optical-audit}
We supply explicit matrix details for the optical identity and isolate
the one-sided endpoint statement used in the trace theorem. These
calculations take place on the original finite itinerary. They do
not assert a bound for a sum over itineraries.

\begin{lemma}[The signed corner cofactor]
\label{lem:v110-cofactor}
For $m\ge1$, with the empty determinant convention when $m=1$,
\begin{equation}\label{eq:v110-signed-corner}
 G_{0m}=(-1)^{m-1}
       \frac{R^2c_0c_m}{(\prod_{i=0}^{m-1}\tau_i)\det J}.
\end{equation}
In particular the absolute value, not the signed mixed entry, gives
$\omega=|G_{0m}|/R$.
\end{lemma}
\begin{proof}
Suppose first that $m>1$ and put $n=m-1$. The minor obtained by
deleting the last row and first column of the tridiagonal matrix
$J$ is triangular, with diagonal entries
$\tau_1^{-1},\ldots,\tau_{n-1}^{-1}$. Cramer's rule therefore gives
\[
 (J^{-1})_{1n}=(-1)^{n+1}
                  \frac{\prod_{i=1}^{n-1}\tau_i^{-1}}{\det J}.
\]
For $n=1$ its numerator is the empty product and the formula says
$(J^{-1})_{11}=1/J_{11}$. The only endpoint--interior couplings are
$H_{01}=R^2c_0c_1/\tau_0$ and
$H_{nm}=R^2c_nc_m/\tau_{m-1}$, whereas $H_{0m}=0$.
Since $H_{II}^{-1}=R^{-2}C_I^{-1}J^{-1}C_I^{-1}$, the Schur
complement gives
\[
 G_{0m}=-\frac{R^2c_0c_m}{\tau_0\tau_{m-1}}(J^{-1})_{1n},
\]
which proves \eqref{eq:v110-signed-corner}. If $m=1$, the mixed
entry follows directly from the single square in
\eqref{eq:v109-second-variation} and has positive sign.

For clarity, the equality between the determinant denominator and
$\mathfrak b_m$ can also be checked without a sign convention for
optical coordinates. Both quantities have initial values
$B_0=0$, $B_1=\tau_0$ and satisfy
\[
 B_{l+1}=\left(1+\tau_l\kappa_l+
                      \frac{\tau_l}{\tau_{l-1}}\right)B_l
                         -\frac{\tau_l}{\tau_{l-1}}B_{l-1}.
\]
For the determinant this is the tridiagonal recurrence after
multiplication by the flight lengths. For the optical product it
follows by eliminating the lower entry of its second column.
This verifies the denominator identity used in
Lemma~\ref{lem:v109-denominator}.
\end{proof}

Put $T_{a,b}=\sum_{i=a}^{b-1}\tau_i$ for $0\le a<b\le m$.
For $S=\{j_1<\cdots<j_s\}\subset\{1,\ldots,m-1\}$ set
$j_0=0$ and $j_{s+1}=m$.

\begin{proposition}[Positive expansion and simultaneous optical marks]
\label{prop:v110-multimark}
The exact denominator has the positive expansion
\begin{equation}\label{eq:v110-positive-polynomial}
 \mathfrak b_m=
 \sum_{S\subset\{1,\ldots,m-1\}}
       \left(\prod_{j\in S}\kappa_j\right)
       \prod_{\nu=0}^{|S|}T_{j_\nu,j_{\nu+1}}.
\end{equation}
For every fixed set $S$ of $s$ distinct interior marks,
\begin{equation}\label{eq:v110-multimark}
 \frac{\omega}{\prod_{j\in S}c_j}
 \le \frac{R_+^{s+1}}{2^s\tau_-^{s+1}}
                       \varrho^{m-s-1}.
\end{equation}
For $S=\varnothing$ this is the first optical bound in
\eqref{eq:v109-optical-bounds}; for $|S|=1$ it is the second.
\end{proposition}
\begin{proof}
Let $E=\left(\begin{smallmatrix}0&0\\1&0\end{smallmatrix}\right)$.
Expand every factor $\mathsf C(\kappa)=I+\kappa E$ in the product
\eqref{eq:v109-optical-product}. Between two selected marks the
shears combine according to
$\mathsf S(a)\mathsf S(b)=\mathsf S(a+b)$, and
$E\mathsf S(T)E=TE$. The top-right entry of the term with marked
set $S$ is exactly the corresponding summand in
\eqref{eq:v110-positive-polynomial}. The empty term is
$T_{0,m}$. Every coefficient is positive.

To obtain the stronger estimate retaining the unmarked reflections,
expand only the factors at the prescribed marks. The coefficient of
$\prod_{j\in S}\kappa_j$, with all other factors retained, is the
product of the top-right entries of the $s+1$ intervening optical
blocks. Indeed $E=e_2e_1^{\mathsf T}$, so repeated multiplication
factors that coefficient into these entries. The block with
$j_{\nu+1}-j_\nu$ flights has top-right entry at least
$\tau_-\varrho^{-(j_{\nu+1}-j_\nu-1)}$, by the same positive-column
argument as in \eqref{eq:v109-denominator-lower}.
Keeping only this nonnegative term gives
\[
 \mathfrak b_m\ge
 \left(\prod_{j\in S}\frac{2}{Rc_j}\right)
            \tau_-^{s+1}\varrho^{-(m-s-1)}.
\]
Substitute this in $\omega=Rc_0c_m/\mathfrak b_m$, and use
$c_0c_m\le1$ and $R\le R_+$, to prove
\eqref{eq:v110-multimark}.
\end{proof}
This estimate controls simultaneous vanishing incidences on a fixed
finite trajectory. Its constant depends on the number of marks,
not on their locations. No word-count factor has been suppressed,
and no pressure conclusion is drawn by summing its right side.

\begin{lemma}[One-sided limits at simultaneous physical mechanisms]
\label{lem:v110-simultaneous}
Fix a positive-length pairing segment and approach one of its ends
from the regular interior. The collision points, velocities, flight
lengths, visited incidences and third-disk clearances have limits
along the entire approach, even if several of the following occur
simultaneously: visited grazing, unvisited tangency, a section
boundary and a critical incoming endpoint. The limiting reduced
Hessian and coarea factor are those in
\eqref{eq:v109-continuous-schur} and \eqref{eq:v109-chi}.
No continuation of the exact event through its boundary is asserted.
\end{lemma}
\begin{proof}
Lemma~\ref{lem:v109-gauss} gives limits for $q$ and $z$ along the
entire segment, not only a subsequence. Hence the initial collision
point and $p_0=-z\cos\theta/R$ converge. The section cap gives the
outgoing normal component $c_0>0$ at the limit, so the initial
velocity converges as well.

Suppose inductively that the departure point $x$ and unit velocity
$v$ of a designated flight converge. For its designated next disk
with centre $O$, the entry time is
\[
 \tau=-(x-O)\cdot v
       -\sqrt{((x-O)\cdot v)^2-|x-O|^2+R^2}.
\]
This is the entry root on the actual regular itinerary. Its
radicand is nonnegative along the approach and has a nonnegative
limit. The root is continuous also when the radicand tends to zero;
$\tau\ge\tau_->0$ prevents a zero-length flight. The arrival point
therefore converges. Reflection by
$v\mapsto v-2(v\cdot n)n$ is continuous even at tangency, so the
next velocity converges. Induction applies a finite number of times.
An earlier disk may become tangent only at the limiting flight;
this does not alter the designated entry-root limit taken from the
original event.

At fixed itinerary and radius, the bounded flight lengths restrict
all possibly nearby third-disk centres to a finite set. Distance
from a compact line segment to each such disk is continuous in the
endpoints. Their finite minimum, and hence the original clearance,
has a limit. A section decision may change at the limiting state,
but its coordinate limit is still obtained from the same side.
None of these continuity arguments requires the boundary mechanisms
to occur separately. Formula \eqref{eq:v109-continuous-schur}, with
$Z^{-1}$ uniformly bounded, then gives the Hessian limit even with
several vanishing incidences. The fixed-direction coarea formula
gives the asserted limit for $\chi$.
\end{proof}

\begin{corollary}[An explicit uniform trace modulus]
\label{cor:v110-holder-constant}
A permissible constant in \eqref{eq:v109-holder} is
\begin{equation}\label{eq:v110-holder-constant}
 C_{\rm H}=
 \frac{500C_{\rm reg}\sqrt{2\Lambda_\circ}}
                 {4\pi c_*\lambda_\circ^2}.
\end{equation}
It is common to all original sheets, collision counts, radii and
exact labels. Its inherited dependence is precisely the regular
variation constant $C_{\rm reg}$; the remaining constants are the
fixed geometry and section constants displayed in this article.
\end{corollary}
\begin{proof}
Corollary~\ref{cor:v109-marked} and $\chi\ge\lambda_\circ$ give
$AI_m\le500\varrho^{(m-2)_+}/(4\pi c_*\lambda_\circ)$.
Integrate the regular derivative bound between $u$ and $v$, and
apply Lemma~\ref{lem:v109-gauss} to that subsegment:
\[
 \int_u^v z^{-1}\,ds
          \le\lambda_\circ^{-1}\sqrt{2\Lambda_\circ(v-u)}.
\]
Multiply the constants and pass to the already proved endpoints.
\end{proof}
Any nonempty collection of visited endpoint grazings belongs to
$\Gamma$, and one vanishing marked incidence is already sufficient
for its zero source statement. Proposition~\ref{prop:v110-multimark}
additionally controls their simultaneous product. This does not
alter the distinct treatment of an unvisited third-disk tangency.
